\chapter{Kết luận và hướng phát triển}
\label{chap:ket_luan}

\section{Kết quả so với mục tiêu}
\label{sec:ket_qua}
\begin{itemize}
    \item \textbf{Lưu trữ và tiền xử lý phân tán}: pipeline \gls{etl} chuyển toàn bộ \num{32000204} bản ghi sang Parquet phân hoạch theo năm bằng PySpark. \dy{(Lộc) Bổ sung: thời gian ETL và tỉ lệ nén thực tế.}
    \item \textbf{Lọc cộng tác và cold-start}: mô hình \gls{als} có checkpoint, tinh chỉnh siêu tham số trên tập kiểm định, sinh Top-20 phim chưa xem cho mỗi người dùng; người dùng mới nhận gợi ý theo Bayesian Popularity. \dy{(Tâm) Bổ sung RMSE và NDCG@10 so với baseline từ E2.}
    \item \textbf{Đánh giá hiệu năng}: \dy{(Lộc) Tóm tắt E3 (mức tuyến tính của Data Scaling) và E4 (Speedup, Efficiency ở 4 nhân, tỉ lệ tuần tự ước lượng).}
\end{itemize}

\section{Hạn chế}
\label{sec:han_che}
\begin{itemize}
    \item Thực nghiệm chạy trên một máy ở chế độ \texttt{local}, chưa đo trên cụm nhiều nút.
    \item Mô hình chỉ dùng điểm đánh giá tường minh; chưa khai thác \texttt{tags.csv}, thể loại hay thời gian.
    \item Gợi ý được tính theo lô; người dùng mới chỉ nhận danh sách phổ biến chung, chưa cá nhân hóa theo vài lượt đánh giá đầu tiên.
    \item Tập kiểm thử chia ngẫu nhiên thay vì theo thời gian.
\end{itemize}

\section{Hướng phát triển}
\label{sec:huong_phat_trien}
\begin{itemize}
    \item Chạy trên cụm Spark nhiều máy (Spark Standalone, YARN hoặc Databricks) để đo Speedup có chi phí mạng.
    \item Dùng ALS phản hồi ngầm (\texttt{implicitPrefs}) hoặc kết hợp thể loại, thẻ để xây dựng mô hình lai.
    \item Gợi ý cho người dùng mới theo cụm phim họ chọn (dựa trên kết quả K-Means), hoặc fold-in vector người dùng từ vài lượt đánh giá đầu tiên.
    \item Chia tập kiểm thử theo thời gian và đánh giá thêm độ đa dạng, độ phủ danh mục của gợi ý.
\end{itemize}
